\documentclass[a4paper,11pt]{article}
\usepackage{../../exercise-style-cmbright}

\renewcommand{\BookTitle}{Precalculus 2e (OpenStax, 2021)}
\renewcommand{\ExerciseSource}{Section}
\renewcommand{\SectionLabel}{1.1}
\renewcommand{\ExerciseNumber}{2}
\renewcommand{\AuthorName}{Henrik Samuelsson}
\renewcommand{\ProblemTitle}{Exercise Solution}

\begin{document}
\maketitle

\begin{problem}
What is the difference between the \emph{input} and the \emph{output} of a function?
\end{problem}

\begin{solution}
The \emph{input} is a value supplied to a function. The \emph{output} is the value the function produces from that input.

A function pairs each input with exactly one output. In each ordered pair, the input is the first value and the output is the second value.

For example, consider a function that pairs each person with their date of birth. The input values could be the names Leonhard Euler and Emmy Noether. The corresponding output values would be 1707-04-15 and 1882-03-23, respectively.
% Birth-date sources:
% https://mathshistory.st-andrews.ac.uk/Biographies/Euler/
% https://mathshistory.st-andrews.ac.uk/Biographies/Noether_Emmy/
\end{solution}

\begin{answer}
The \emph{input} is a value supplied to a function. The \emph{output} is the value the function produces from that input.
\end{answer}

\end{document}
